%% example.tex — Beispielkapitel mit ALLEN eigenen Umgebungen
%% Demonstriert: process, infobox, warnbox, procon, codebox,
%%               \sidetables, \stdtable, Abbildung, \cite{}.
%% ──────────────────────────────────────────────────────────────────────────

\section{Grundlagen der Algorithmen}

\subsection{Definition}

\begin{infobox}[Definition: Algorithmus]
Ein \textbf{Algorithmus} ist eine endliche, eindeutige Folge von Anweisungen,
die ein Problem in einer endlichen Anzahl von Schritten löst
\parencite{cormen2022}.
Er besitzt stets eine klare Eingabe, eine definierte Ausgabe und terminiert
nach endlich vielen Schritten.
\end{infobox}

Algorithmen sind das Herzstück der Informatik. Seit \textcite{dijkstra1959}
seinen wegweisenden Kürzeste-Wege-Algorithmus vorstellte, ist die formale
Analyse von Algorithmen ein eigenständiges Forschungsfeld.

\subsection{Eigenschaften eines korrekten Algorithmus}

Ein Algorithmus muss folgende Eigenschaften erfüllen \parencite{cormen2022}:

\begin{process}
  \step{%
    \textbf{Finitheit:} Der Algorithmus endet nach einer endlichen Anzahl von
    Schritten unter allen zulässigen Eingaben.%
  }
  \step{%
    \textbf{Eindeutigkeit:} Jede Anweisung ist klar und unmissverständlich
    definiert; kein Schritt lässt Interpretationsspielraum.%
  }
  \step{%
    \textbf{Ausführbarkeit:} Alle Schritte sind mit den zur Verfügung stehenden
    Mitteln tatsächlich durchführbar.%
  }
  \step{%
    \textbf{Determiniertheit:} Gleiche Eingaben liefern stets dasselbe Ergebnis
    (bei deterministischen Algorithmen).%
  }
  \step{%
    \textbf{Effizienz:} Der Ressourcenverbrauch (Zeit, Speicher) ist in einem
    akzeptablen Verhältnis zur Problemgröße.%
  }
\end{process}

\subsection{Vorteile und Nachteile formaler Algorithmen}

\begin{procon}
  \pro{Wiederverwendbar und übertragbar auf verschiedene Probleminstanzen}
  \pro{Formale Korrektheitsnachweise möglich}
  \pro{Laufzeit- und Speicherverbrauch systematisch analysierbar}
  \pro{Unabhängig von Programmiersprache oder Plattform beschreibbar}
  \con{Formale Beschreibung kann komplex und schwer lesbar werden}
  \con{Optimierung für konkrete Hardware oft notwendig}
  \con{Nicht alle Probleme sind algorithmisch lösbar (Halteproblem)}
  \con{Erheblicher Aufwand für den Korrektheitsbeweis}
\end{procon}

\subsection{Wichtiger Hinweis zur Komplexitätsanalyse}

\begin{warnbox}[Achtung: Asymptotische Notation]
Die \textbf{O-Notation} beschreibt das \emph{asymptotische} Wachstum und
ignoriert konstante Faktoren sowie untere Terme.
Ein Algorithmus mit $O(n \log n)$ kann für kleine $n$ langsamer sein als
ein $O(n^2)$-Algorithmus mit kleinen Konstanten!
Stets auch den \emph{praktischen} Anwendungsfall berücksichtigen.
\end{warnbox}

% ──────────────────────────────────────────────────────────────────────────
\section{Laufzeitkomplexität}

\subsection{Überblick und Notation}

\begin{infobox}[Wichtige O-Klassen (von schnell nach langsam)]
\begin{tabular}{@{}ll@{}}
  $O(1)$        & konstant           \\
  $O(\log n)$   & logarithmisch      \\
  $O(n)$        & linear             \\
  $O(n \log n)$ & quasilinear        \\
  $O(n^2)$      & quadratisch        \\
  $O(2^n)$      & exponentiell       \\
  $O(n!)$       & faktoriell         \\
\end{tabular}
\end{infobox}

\subsection{Abbildung: Wachstum der O-Klassen}

\begin{figure}[H]
  \centering
  % Platzhalter – ersetze durch: \includegraphics[width=0.7\linewidth]{komplexitaet}
  \fbox{\parbox{0.65\linewidth}{\centering\vspace{2cm}
    \color{primary!50!white}\Large\sffamily
    Abbildung: assets/images/komplexitaet.pdf\\[4pt]
    \small (Platzhaltergrafik – Datei in assets/images/ ablegen)
    \vspace{2cm}}}
  \caption{Wachstumsverhalten verschiedener Komplexitätsklassen für
           $n \in [1, 20]$. Quelle: eigene Darstellung nach
           \textcite{cormen2022}.}
  \label{fig:komplexitaet}
\end{figure}

\subsection{Vergleich bekannter Sortieralgorithmen}

Die folgende Tabelle vergleicht gebräuchliche Sortierverfahren anhand ihrer
asymptotischen Laufzeiten (vgl.\ \textcite{cormen2022}).

\begin{stdtable}{l X X X l}{Algorithmus & Best Case & Average Case & Worst Case & Stabil?}
  Bubble Sort    & $O(n)$         & $O(n^2)$       & $O(n^2)$       & Ja  \\
  Selection Sort & $O(n^2)$       & $O(n^2)$       & $O(n^2)$       & Nein\\
  Insertion Sort & $O(n)$         & $O(n^2)$       & $O(n^2)$       & Ja  \\
  Merge Sort     & $O(n \log n)$  & $O(n \log n)$  & $O(n \log n)$  & Ja  \\
  Quick Sort     & $O(n \log n)$  & $O(n \log n)$  & $O(n^2)$       & Nein\\
  Heap Sort      & $O(n \log n)$  & $O(n \log n)$  & $O(n \log n)$  & Nein\\
  Counting Sort  & $O(n+k)$       & $O(n+k)$       & $O(n+k)$       & Ja  \\
\end{stdtable}

\subsubsection{Zwei Tabellen nebeneinander (\texttt{\textbackslash sidetables})}

\sidetables{%
  \textbf{Interne Sortierverfahren}\medskip

  \begin{tabular}{@{}lc@{}}
    \toprule
    Verfahren      & In-Place? \\
    \midrule
    Insertion Sort & Ja        \\
    Quick Sort     & Ja        \\
    Merge Sort     & Nein      \\
    Heap Sort      & Ja        \\
    \bottomrule
  \end{tabular}%
}{%
  \textbf{Externe Sortierverfahren}\medskip

  \begin{tabular}{@{}lc@{}}
    \toprule
    Verfahren         & Passes \\
    \midrule
    Externes Merge    & $O(\log n)$ \\
    Polyphasen-Merge  & $O(\log n)$ \\
    Replacement Sel.  & $O(n)$      \\
    \bottomrule
  \end{tabular}%
}

% ──────────────────────────────────────────────────────────────────────────
\section{Implementierungsbeispiel}

\subsection{Merge Sort in Python}

Der folgende Code zeigt eine kompakte Python-Implementierung von Merge Sort.
Eine ausführliche Erklärung findet sich bei \textcite{wikipedia2024}.

\begin{codebox}
def merge_sort(arr):
    if len(arr) <= 1:
        return arr

    mid   = len(arr) // 2
    left  = merge_sort(arr[:mid])
    right = merge_sort(arr[mid:])
    return merge(left, right)

def merge(left, right):
    result, i, j = [], 0, 0
    while i < len(left) and j < len(right):
        if left[i] <= right[j]:
            result.append(left[i]); i += 1
        else:
            result.append(right[j]); j += 1
    return result + left[i:] + right[j:]

# Beispielaufruf
data = [38, 27, 43, 3, 9, 82, 10]
print(merge_sort(data))   # -> [3, 9, 10, 27, 38, 43, 82]
\end{codebox}

\subsection{Laufzeit-Analyse des Merge-Sort-Schritts}

\begin{process}
  \step{Teile das Array der Länge $n$ in zwei Hälften der Länge $\lfloor n/2 \rfloor$
        und $\lceil n/2 \rceil$.}
  \step{Sortiere jede Hälfte rekursiv — Laufzeit je $T(n/2)$.}
  \step{Füge die zwei sortierten Hälften zusammen — lineare Arbeit $O(n)$.}
  \step{Die Rekurrenz $T(n) = 2\,T(n/2) + O(n)$ ergibt nach dem Master-Theorem
        $T(n) = O(n \log n)$.}
\end{process}

\begin{infobox}[Master-Theorem (Fall 2)]
Gilt $T(n) = a\,T(n/b) + \Theta(n^{\log_b a})$, so ist
$T(n) = \Theta(n^{\log_b a} \log n)$.
Für Merge Sort: $a = 2,\; b = 2,\; f(n) = n = \Theta(n^{\log_2 2}) = \Theta(n)$,
also $T(n) = \Theta(n \log n)$.
\end{infobox}
